% chapters/15/praktikum.tex
\begin{praktikumbox}
Praktikum ini mengikuti Tugas Praktikum 12 pada RPS: menambahkan
\code{AudioSource} dan \code{AudioClip} untuk musik latar serta efek suara,
memicu efek suara pada peristiwa lompat, mengambil benda, dan permainan
berakhir, serta menyesuaikan \istilah{timing} dan tingkat kekuatannya.

\begin{langkah}
  \item \textbf{Kumpulkan berkas bunyinya.} Sediakan satu berkas musik latar
        dan sekurang-kurangnya empat berkas efek suara: lompat, mengambil
        benda, terluka, dan kalah. Catat sumber serta lisensi setiap berkas
        sejak awal.

  \item \textbf{Impor dan atur.} Seret berkasnya ke folder \code{Assets}. Pada
        berkas musik, matikan \istilah{Load In Background} apabila musiknya
        perlu mulai seketika; pada efek suara pendek, pilih
        \istilah{Decompress On Load} agar tidak ada penundaan.

  \item \textbf{Susun pemusat bunyinya.} Buat \istilah{GameObject} kosong
        bernama \code{ManajerBunyi}, lalu tambahkan dua komponen
        \code{AudioSource} ke dalamnya: satu bernama \code{musicSource} dan
        satu bernama \code{sfxSource}. Tempelkan skrip
        \code{ch15_manajer_bunyi.cs} dari Daftar~\ref{lst:ch15-manajer}, lalu
        hubungkan keenam rujukannya pada \istilah{Inspector}.

  \item \textbf{Atur ruang bunyinya.} Pada \code{musicSource}, atur
        \istilah{Spatial Blend} menjadi nol, sebab musik latar tidak memiliki
        tempat. Pada \code{sfxSource}, nilai itu juga dibuat nol untuk efek
        suara yang selalu berkaitan dengan pemain.

  \item \textbf{Tambahkan pemicunya pada pemain.} Tempelkan skrip
        \code{ch15_pemicu_bunyi.cs} dari Daftar~\ref{lst:ch15-pemicu} pada
        benda pemain, lalu hubungkan rujukan \code{manajer} dan \code{hud}.
        Tambahkan dua baris pada \code{ch13_kontrol_pemain.cs} sebagaimana
        diperlihatkan pada Bagian~\ref{sec:peristiwa-bunyi}, agar bunyi lompat
        terpanggil.

  \item \textbf{Tandai benda yang dapat diambil.} Buat tanda \code{Pickup}
        melalui menu \istilah{Tag}, lalu pasang pada setiap benda yang dapat
        diambil. Tanpa langkah ini, sentuhannya diabaikan tanpa keluhan apa
        pun.

  \item \textbf{Susun tingkat kekuatannya.} Buat \istilah{AudioMixer} melalui
        menu \istilah{Window}, lalu \istilah{Audio}, lalu
        \istilah{Audio Mixer}. Tambahkan dua tingkat di bawah \code{Master},
        yaitu \code{Musik} dan \code{Efek}, lalu arahkan keluaran setiap
        \code{AudioSource} ke tingkat yang sesuai.

  \item \textbf{Jalankan empat percobaan} pada
        Tabel~\ref{tab:percobaan-bunyi} satu per satu. Catat apa yang Anda
        dengar pada setiap percobaan, sebelum dan sesudah nilainya diubah.

  \item \textbf{Tambahkan partikel sederhana} pada satu peristiwa, misalnya
        pada saat benda diambil, sebagai pasangan yang terlihat bagi bunyinya.

  \item \textbf{Tulis catatan praktikum} yang memuat daftar berkas bunyi
        beserta sumber dan lisensinya, daftar peristiwa beserta bunyinya,
        nilai tingkat kekuatan yang dipakai, dan hasil keempat percobaan.

  \item \textbf{Simpan pekerjaan} ke repositori Git. Berkas bunyi berukuran
        besar, sehingga berkasnya sebaiknya disimpan di luar repositori dan
        alamatnya dicatat pada berkas keterangan.
\end{langkah}
\end{praktikumbox}

\begin{table}[htbp]
\centering
\caption{Lembar kerja Praktikum 12}
\label{tab:lembar-praktikum-12}
{\setlength{\tabcolsep}{4pt}\footnotesize%
\begin{tabularx}{\textwidth}{@{}L{3.6cm} L{3.4cm} Y@{}}
\toprule
\headrow \textbf{Peristiwa} & \textbf{Berkas bunyi} & \textbf{Tingkat kekuatan dan catatan} \\
\midrule
Scene dimulai & & \\[1.4em]
\midrule
Pemain melompat & & \\[1.4em]
\midrule
Mengambil benda & & \\[1.4em]
\midrule
Pemain terluka & & \\[1.4em]
\midrule
Permainan berakhir & & \\[1.4em]
\midrule
\addlinespace
\multicolumn{3}{@{}p{\textwidth}@{}}{\footnotesize\itshape
Berkas musik latar beserta sumber dan lisensinya: \dotfill} \\[1.4em]
\addlinespace
\multicolumn{3}{@{}p{\textwidth}@{}}{\footnotesize\itshape
Hasil percobaan terpenting yang Anda lakukan: \dotfill} \\[1.8em]
\bottomrule
\end{tabularx}}
\end{table}

\begin{catatabox}
Langkah kedelapan adalah langkah yang paling penting bagi indikator \ind{c}.
Mengubah satu nilai lalu mendengarkan akibatnya adalah cara satu-satunya untuk
memahami tingkat kekuatan bunyi. Nilai yang tercantum pada bab ini adalah
titik awal; yang menentukan adalah pendengaran Anda pada level Anda sendiri.
\end{catatabox}
